% GENERATED FILE. Edit canonical lessons, then run npm run generate:books.
% canonical-source-hash: fnv1a64:98a23b6385f63ae4
% canonical-lessons: FR-C249-acteur, FR-C249-actrice, FR-C249-poeme, FR-C249-ecrivain, FR-C249-peintre

\chapter{Actor, Actress, Poem, Writer, Painter}
\label{ch:fr-actor-actress-poem-writer-painter}
\hlchaptermodality{249}

\begin{chapteropening}
\textbf{By the end of this chapter:} \emph{I can say an actor, an actress, a poem, a writer and a painter.}

Complete the last lesson of chapter 249: I can say an actor, an actress, a poem, a writer and a painter.
\end{chapteropening}

\section[l'acteur]{l'acteur — an actor}

\label{lesson:FR-C249-acteur}



\begin{quote}
Before the new one: say the French for politics, then the French for the economy.
\end{quote}

\subsection*{You'll want to know: l'acteur}
\textbf{l'acteur} — \textquotedblleft{}an actor\textquotedblright{}.

\begin{grammarlens}[title={What you've built}]
One more word for this chapter.
\end{grammarlens}

\subsection*{Your turn}
\begin{itemize}
\raggedright
  \item \emph{Say it:} \emph{l'acteur}
  \item \emph{Say it:} \emph{l'acteur}, once more
  \item \emph{Say it:} say \emph{l'acteur}
  \item \emph{Recall:} say the French for politics, then the French for the economy, then say \emph{l'acteur} again
\end{itemize}

\begin{tcolorbox}[breakable,colback=teal!4,colframe=teal!35!black,title={Before you move on}]
What does l'acteur mean? (\textquotedblleft{}An actor\textquotedblright{}.) Say it once more.
\end{tcolorbox}
\section[l'actrice]{l'actrice — an actress}

\label{lesson:FR-C249-actrice}



\begin{quote}
Before the new one: say the French for the economy, then the French for an actor.
\end{quote}

\subsection*{You'll want to know: l'actrice}
\textbf{l'actrice} — \textquotedblleft{}an actress\textquotedblright{}.

\begin{grammarlens}[title={What you've built}]
One more word for this chapter.
\end{grammarlens}

\subsection*{Your turn}
\begin{itemize}
\raggedright
  \item \emph{Say it:} \emph{l'actrice}
  \item \emph{Say it:} \emph{l'actrice}, once more
  \item \emph{Say it:} say \emph{l'actrice}
  \item \emph{Recall:} say the French for the economy, then the French for an actor, then say \emph{l'actrice} again
\end{itemize}

\begin{tcolorbox}[breakable,colback=teal!4,colframe=teal!35!black,title={Before you move on}]
What does l'actrice mean? (\textquotedblleft{}An actress\textquotedblright{}.) Say it once more.
\end{tcolorbox}
\section[le poème]{le poème — a poem}

\label{lesson:FR-C249-poeme}



\begin{quote}
Before the new one: say the French for an actor, then the French for an actress.
\end{quote}

\subsection*{You'll want to know: le poème}
\textbf{le poème} — \textquotedblleft{}a poem\textquotedblright{}.

\begin{grammarlens}[title={What you've built}]
One more word for this chapter.
\end{grammarlens}

\subsection*{Your turn}
\begin{itemize}
\raggedright
  \item \emph{Say it:} \emph{le poème}
  \item \emph{Say it:} \emph{le poème}, once more
  \item \emph{Say it:} say \emph{le poème}
  \item \emph{Recall:} say the French for an actor, then the French for an actress, then say \emph{le poème} again
\end{itemize}

\begin{tcolorbox}[breakable,colback=teal!4,colframe=teal!35!black,title={Before you move on}]
What does le poème mean? (\textquotedblleft{}A poem\textquotedblright{}.) Say it once more.
\end{tcolorbox}
\section[l'écrivain]{l'écrivain — a writer}

\label{lesson:FR-C249-ecrivain}



\begin{quote}
Before the new one: say the French for an actress, then the French for a poem.
\end{quote}

\subsection*{You'll want to know: l'écrivain}
\textbf{l'écrivain} — \textquotedblleft{}a writer\textquotedblright{}.

\begin{grammarlens}[title={What you've built}]
One more word for this chapter.
\end{grammarlens}

\subsection*{Your turn}
\begin{itemize}
\raggedright
  \item \emph{Say it:} \emph{l'écrivain}
  \item \emph{Say it:} \emph{l'écrivain}, once more
  \item \emph{Say it:} say \emph{l'écrivain}
  \item \emph{Recall:} say the French for an actress, then the French for a poem, then say \emph{l'écrivain} again
\end{itemize}

\begin{tcolorbox}[breakable,colback=teal!4,colframe=teal!35!black,title={Before you move on}]
What does l'écrivain mean? (\textquotedblleft{}A writer\textquotedblright{}.) Say it once more.
\end{tcolorbox}
\section[le peintre]{le peintre — a painter}

\label{lesson:FR-C249-peintre}



\begin{quote}
Before the new one: say the French for a poem, then the French for a writer.
\end{quote}

\subsection*{You'll want to know: le peintre}
\textbf{le peintre} — \textquotedblleft{}a painter\textquotedblright{}.

\begin{grammarlens}[title={What you've built}]
One more word for this chapter.
\end{grammarlens}

\subsection*{Your turn}
\begin{itemize}
\raggedright
  \item \emph{Say it:} \emph{le peintre}
  \item \emph{Say it:} \emph{le peintre}, once more
  \item \emph{Say it:} say \emph{le peintre}
  \item \emph{Recall:} say the French for a poem, then the French for a writer, then say \emph{le peintre} again
\end{itemize}

\begin{tcolorbox}[breakable,colback=teal!4,colframe=teal!35!black,title={Before you move on}]
What does le peintre mean? (\textquotedblleft{}A painter\textquotedblright{}.) Say it once more.
\end{tcolorbox}
